\documentclass[11pt]{article}
\usepackage{amsmath,amssymb,amsthm}
\usepackage[margin=1in]{geometry}
\newtheorem{theorem}{Theorem}
\newtheorem{lemma}{Lemma}
\newtheorem{proposition}{Proposition}
\newtheorem{corollary}{Corollary}
\newtheorem{conjecture}{Conjecture}
\theoremstyle{definition}
\newtheorem{remark}{Remark}
\DeclareMathOperator{\val}{val}
\DeclareMathOperator{\Res}{Res}
\newcommand{\la}{\lambda}
\newcommand{\C}{\mathbb{C}}
\newcommand{\Z}{\mathbb{Z}}
\newcommand{\sll}{s_\la}
\newcommand{\inner}[2]{\langle #1,#2\rangle}

\title{The $d=4$ fibre--vanishing law $G_\la(i)=0\iff\la=(2,2)$:\\
structural reductions, the two--row case, and an isolated gap}
\author{Clio}
\date{2026-06-04}

\begin{document}
\maketitle

\begin{abstract}
For $\la\vdash n$ let $G_\la(q)=\sum_{T\in\mathrm{SYT}(\la)}q^{s(T)}$ be the graded
descent enumerator with the parity--twisted statistic $s(T)=\sum_{i\in\mathrm{Des}(T)}w_i$,
$w_i=2i-1$ if $n-i$ is odd and $0$ otherwise.  We study the fibre at $\zeta_4=i$.
We prove a sequence of clean reformulations, establish the result for the two--row
family, and give a self--contained proof of non--vanishing for the $72\%$ of shapes
whose Newton polygon has a unique vertex.  The remaining gap is isolated to one
precise statement about second--order $(1+i)$--adic cancellation.  \textbf{Status:
partial.}  The full biconditional is verified computationally for all $\la$ with
$n\le 16$, with $(2,2)$ the unique vanisher.
\end{abstract}

\section{Setup and the master form}
Throughout $n=2m$ is even.  (For $n=2m+1$ odd the master identity gives
$G_\la(i)=\inner{\sll}{h_1(h_2-ie_2)^m}$; computation shows no vanishers for $n\le16$,
and since the lone vanisher $(2,2)$ has $n=4$ even, the odd case is not needed for
Theorem~\ref{thm:target} as stated.)  The proven master
identity specialised at $\zeta=i$ gives the \emph{Gaussian--integer form}
\begin{equation}\label{eq:gauss}
G_\la(i)=\inner{\sll}{\psi^m},\qquad \psi:=h_2+i\,e_2=s_2+i\,s_{11},
\end{equation}
where $\inner{\cdot}{\cdot}$ is the Hall inner product.  Equivalently, with
\begin{equation}\label{eq:P}
P_\la(t):=\inner{\sll}{(h_2+t\,e_2)^m}=\sum_{v=0}^m c_v\,t^v,
\qquad c_v=\inner{\sll}{h_2^{\,m-v}e_2^{\,v}}\in\Z_{\ge0},
\end{equation}
we have $G_\la(i)=P_\la(i)$.  The $c_v$ are nonnegative because they count
$2$--strip chains: $c_v$ is the number of saturated chains
$\varnothing=\mu^0\subset\mu^1\subset\dots\subset\mu^m=\la$ with each step a
$2$--box horizontal or vertical strip, exactly $v$ of them vertical (a disconnected
$2$--box step counts in both classes, via the Pieri coefficients
$\inner{s_{\mu^j}}{s_2 s_{\mu^{j-1}}}$ and $\inner{s_{\mu^j}}{s_{11}s_{\mu^{j-1}}}$).
Basic invariants: $P_\la(1)=f^\la$, $\deg P_\la=m$, $P_\la(-1)=\chi^\la(2^m)$.

\begin{theorem}[Target]\label{thm:target}
For $\la\vdash 2m$, $\;G_\la(i)=0\iff\la=(2,2)$.
\end{theorem}

Because $P_\la$ has real coefficients and $\pm i$ are conjugate, $i$ is a root iff
$-i$ is, so Theorem~\ref{thm:target} is equivalent to each of:
\begin{itemize}
\item $(t^2+1)\mid P_\la(t)$ only for $\la=(2,2)$;
\item $R_\la:=\sum_{v\ \mathrm{even}}(-1)^{v/2}c_v=0$ and
      $I_\la:=\sum_{v\ \mathrm{odd}}(-1)^{(v-1)/2}c_v=0$ simultaneously only for $(2,2)$;
\item $|G_\la(i)|^2=\Res(t^2+1,\,P_\la(t))=0$ only for $(2,2)$.
\end{itemize}
The lone solution is clean: $P_{(2,2)}(t)=1+t^2$.

\section{Two structural identities}

\begin{lemma}[Conjugation]\label{lem:conj}
$P_{\la'}(t)=t^m P_\la(1/t)$, i.e.\ $c_v(\la')=c_{m-v}(\la)$.  Hence
$R_{\la'}+iI_{\la'}=i^m\bigl(R_\la-iI_\la\bigr)$, so the loci $\{R=0\}$ and
$\{I=0\}$ are interchanged (for $m$ odd) or each preserved (for $m$ even) under
conjugation; in particular $G_\la(i)=0\iff G_{\la'}(i)=0$.
\end{lemma}
\begin{proof}
Apply the involution $\omega$ ($\omega h_2=e_2$, $\omega e_2=h_2$, $\omega\sll=s_{\la'}$,
an isometry) to \eqref{eq:P}: $P_\la(t)=\inner{s_{\la'}}{(e_2+t h_2)^m}
=t^m\inner{s_{\la'}}{(h_2+t^{-1}e_2)^m}=t^m P_{\la'}(1/t)$.
\end{proof}

\begin{lemma}[Coproduct of $\psi$]\label{lem:cop}
In the Hopf algebra of symmetric functions,
\[
\Delta\psi=\psi\otimes1+1\otimes\psi+(1+i)\,(p_1\otimes p_1).
\]
Consequently the skewing operator $\psi^\perp$ is a second--order operator
$\psi^\perp=\tfrac{1+i}{2}\bigl(\partial_{p_1}^2-2i\,\partial_{p_2}\bigr)$, and
$G_\la(i)=\bigl((\psi^\perp)^m\sll\bigr)\big|_{\varnothing}$.
\end{lemma}
\begin{proof}
$\Delta h_2=h_2\otimes1+h_1\otimes h_1+1\otimes h_2$ and likewise for $e_2$ with
$e_1=h_1=p_1$; adding with weight $i$ gives $h_1\otimes h_1+i\,e_1\otimes e_1
=(1+i)p_1\otimes p_1$.  The operator form follows from
$\psi=\tfrac{1+i}{2}(p_1^2-ip_2)$ and $p_k^\perp=k\,\partial_{p_k}$.
\end{proof}

\noindent These give the down--operator recursion
\begin{equation}\label{eq:rec}
P_\la(i)=\sum_{\mu:\ \la/\mu\ \mathrm{a\ 2\text{-}strip}}\omega(\la/\mu)\,P_\mu(i),
\qquad \omega\in\{1,\ i,\ 1+i\},
\end{equation}
the weight being $1$ for a horizontal domino, $i$ for a vertical domino, and
$1+i$ for a disconnected pair.  Iterating, $G_\la(i)=\sum_{\text{chains}}\prod\omega$.

\section{The multivariate generating function}
The single most useful reformulation: since $\sll(x_1,\dots,x_N)=0$ for
$N<\ell(\la)$ and $\inner{\sll}{f}$ equals the $\sll$--coefficient of $f(x_1,\dots,x_N)$
for $N\ge\ell(\la)$,
\begin{equation}\label{eq:genfun}
G_\la(i)=\bigl[\,\sll\text{--coefficient of}\,\bigr]\ \psi(x_1,\dots,x_N)^m,
\qquad \psi(x)=\sum_i x_i^2+(1+i)\!\sum_{i<j}x_ix_j,
\end{equation}
with $N=\ell(\la)$.  Equivalently, with Vandermonde
$a_\delta=\prod_{i<j}(x_i-x_j)$ and $\delta=(N-1,\dots,1,0)$,
\[
G_\la(i)=\bigl[x^{\la+\delta}\bigr]\,\bigl(a_\delta(x)\,\psi(x)^m\bigr).
\]

\section{The two--row family}
Here $N=2$, $\psi(u,v)=u^2+v^2+(1+i)uv=(u-\rho v)(u-\sigma v)$ where $\rho,\sigma$
are the roots of $z^2+(1+i)z+1=0$ (so $\rho\sigma=1$, $\rho+\sigma=-(1+i)$, and
$|\rho|\approx0.587<1<|\sigma|\approx1.704$; neither lies on the unit circle).

\begin{proposition}\label{prop:tworow}
Write $R(u):=(u^2+(1+i)u+1)^m=\sum_{l=0}^{2m}r_l u^l$, a palindromic polynomial
($r_l=r_{2m-l}$), with
\[
r_l=\sum_{2p+q=l}(1+i)^q\,\frac{m!}{p!\,q!\,(m-p-q)!}.
\]
Then for $a+b=2m$, $a\ge b$, one has
\[
G_{(a,b)}(i)=[u^{a+1}]\bigl((u-1)R(u)\bigr)=r_a-r_{a+1}
=r_b-r_{b-1}\quad(\text{using }r_l=r_{2m-l}).
\]
\end{proposition}
\begin{proof}
By \eqref{eq:genfun} with $N=2$, $\psi(u,v)^m=\sum_{a\ge b}G_{(a,b)}s_{(a,b)}(u,v)$;
multiplying by the Vandermonde $a_\delta=u-v$ and using
$s_{(a,b)}(u,v)\,(u-v)=u^{a+1}v^b-u^bv^{a+1}$ gives
$(u-v)\psi(u,v)^m=\sum_{a\ge b}G_{(a,b)}(u^{a+1}v^b-u^bv^{a+1})$.  Setting $v=1$
and reading the coefficient of $u^{a+1}$ (with $a+1>b$) yields
$G_{(a,b)}=[u^{a+1}]((u-1)R(u))=r_a-r_{a+1}$.  The trinomial expansion of
$R=(u^2+(1+i)u+1)^m$ gives the stated $r_l$.
\end{proof}

\begin{corollary}\label{cor:tworow-low}
For the two--row shapes one has the explicit low--$b$ values
\[
G_{(2m,0)}=1,\quad G_{(2m-1,1)}=(m-1)+mi,\quad G_{(2m-2,2)}=m(m-2)\,i .
\]
In particular among the families $b\in\{0,1,2\}$ the unique vanisher is
$G_{(2,2)}=0$ (at $m=2$).  Moreover $G_{(a,b)}=0\iff r_b=r_{b-1}$, and this is
verified to have $(2,2)$ as its only solution for all $m\le 29$.
\end{corollary}
\begin{proof}
$r_{2m}=1$ gives $G_{(2m,0)}=r_{2m}-r_{2m+1}=1$.  For $b=1$:
$r_{2m-1}=m(1+i)$, $r_{2m}=1$, so $G=m(1+i)-1=(m-1)+mi$, with imaginary part
$m\ge1$.  For $b=2$: $r_{2m-2}=\binom m2(1+i)^2+\binom m1=m(m-1)i+m$ and
$r_{2m-1}=m(1+i)$, so $G=m(m-1)i+m-m(1+i)=m(m-2)i$, zero iff $m=2$.
\end{proof}

\begin{remark}
Corollary~\ref{cor:tworow-low} already exhibits the mechanism that singles out
$(2,2)$: inside the natural sub--family $(2m-2,2)$ the real part vanishes
identically and the imaginary part is $m(m-2)$, which has its only positive root
at $m=2$.  Thus $(2,2)$ is ``the $(2m-2,2)$ shape at the unique $m$ where the
quadratic $m(m-2)$ vanishes.''  A uniform proof that $r_b\ne r_{b-1}$ for all
$0\le b\le m$, $(m,b)\ne(2,2)$, would close the two--row case; the obstruction is
that $r_b-r_{b-1}$ is a $2$--dimensional ($\mathbb Z[i]$) sum whose real and
imaginary parts must be shown not to vanish simultaneously, and these change sign
with $b$.  This remains open (Gap~B below).
\end{remark}

\section{The $(1+i)$--adic Newton polygon: $72\%$ unconditionally}
Substitute $h_1^2=h_2+e_2$ to pass to the basis adapted to $\pi:=1+i$.  Define
\begin{equation}\label{eq:A}
A_\la(x):=\inner{\sll}{(h_1^2+x\,e_2)^m}=P_\la(1+x)=\sum_{j=0}^m a_j x^j,
\qquad a_j=\binom mj M_j,\ \ M_j=\inner{\sll}{h_1^{2(m-j)}e_2^j}\in\Z_{\ge0}.
\end{equation}
Then $G_\la(i)=A_\la(i-1)$ and $i-1=i\pi$ has $\pi$--adic valuation $v_\pi(i-1)=1$.
Work in $\Z[i]$ with the discrete valuation $v_\pi$ (so $v_\pi(2)=2$ and
$v_\pi(N)=2v_2(N)$ for $N\in\Z$).  Set
\[
\val(j):=v_\pi\bigl(a_j(i\pi)^j\bigr)=j+2v_2(a_j)\qquad(a_j\ne0),
\qquad V:=\min_{a_j\ne0}\val(j),\quad J^\star:=\{j:\val(j)=V\}.
\]

\begin{lemma}\label{lem:parity}
$\val(j+1)-\val(j)$ is odd for every $j$ with $a_j,a_{j+1}\ne0$.  Consequently
$J^\star$ lies in a single parity class of $j$ (no two consecutive indices tie).
\end{lemma}
\begin{proof}
$\val(j+1)-\val(j)=1+2\bigl(v_2(a_{j+1})-v_2(a_j)\bigr)$ is $1+(\text{even})$.
\end{proof}

\begin{lemma}[Leading residue]\label{lem:lead}
Write $G_\la(i)=A_\la(i\pi)=\sum_j a_j(i\pi)^j$.  Then
$v_\pi(G_\la)\ge V$, and modulo $\pi$,
\[
\pi^{-V}G_\la(i)\equiv |J^\star|\pmod\pi .
\]
\end{lemma}
\begin{proof}
For $j\in J^\star$ write $a_j=2^{e_j}b_j$ with $b_j$ odd.  Using $2=-i\pi^2$,
\[
a_j(i\pi)^j=2^{e_j}b_j\,i^j\pi^j=\pi^{\,j+2e_j}\,i^j(-i)^{e_j}b_j=\pi^{V}\,u_j,
\quad u_j:=i^j(-i)^{e_j}b_j .
\]
The residue field $\Z[i]/(\pi)\cong\mathbb F_2$ has $i\equiv1$ and $b_j\equiv1$, so
$u_j\equiv1\pmod\pi$.  Terms with $j\notin J^\star$ have $\val(j)\ge V+1$, hence
$v_\pi\ge V+1$.  Summing, $\pi^{-V}G_\la\equiv\sum_{j\in J^\star}1=|J^\star|$.
\end{proof}

\begin{theorem}[Unique--vertex non--vanishing]\label{thm:partial}
If $|J^\star|=1$, or more generally if $|J^\star|$ is odd, then $v_\pi(G_\la)=V<\infty$
and hence $G_\la(i)\ne0$.  This hypothesis holds for $212$ of the $294$ shapes
with $n\le14$ ($72.1\%$), and for every such shape $v_\pi(G_\la)$ equals the
predicted $V$ (verified).
\end{theorem}
\begin{proof}
Immediate from Lemma~\ref{lem:lead}: if $|J^\star|$ is odd then
$\pi^{-V}G_\la\equiv1\not\equiv0\pmod\pi$, so $v_\pi(G_\la)=V$, finite.
\end{proof}

\section{The isolated gaps}
\paragraph{Gap A (the even ties).}  Theorem~\ref{thm:partial} fails exactly when
$|J^\star|$ is even, which empirically is \emph{always} the case at a genuine tie
($|J^\star|\in\{2,4\}$ for all $n\le14$).  Then the leading residue vanishes and
$v_\pi(G_\la)\ge V+1$.  For a size--$2$ tie $J^\star=\{p,p+2\}$ one computes the
next layer explicitly: with $e:=v_2(a_{p+2})$ and $v_2(a_p)=e+1$ (forced by
$\val(p)=\val(p+2)$), writing $a_p=2^{e+1}a_p'$, $a_{p+2}=2^e a_{p+2}'$
($a_p',a_{p+2}'$ odd),
\[
a_p(i\pi)^p+a_{p+2}(i\pi)^{p+2}=(i\pi)^p\bigl(a_p-2i\,a_{p+2}\bigr)
=2^{e+1}(a_p'-i\,a_{p+2}')\,(i\pi)^p,
\]
and $a_p'-i a_{p+2}'$ has $v_\pi=1$ (sum of two odd squares is $\equiv2\bmod8$),
so this contributes at level $V+1$.  But the data show $v_\pi(G_\la)$ is often
$V+2$ or higher, so the level--$(V+1)$ contribution is itself cancelled by the
opposite--parity terms (which Lemma~\ref{lem:parity} places exactly one level up).
\textbf{What is needed:} an $m$--uniform argument that this alternating, multi--level
$\pi$--adic cancellation \emph{terminates} (i.e.\ $v_\pi(G_\la)<\infty$) for every
$\la\ne(2,2)$.  The shape $(2,2)$ is the unique one where it descends forever
($A_{(2,2)}(x)=x^2+2x+2$ is exactly the minimal polynomial of $i-1$).
The cancellation depth is genuinely unbounded: the maximum finite $v_\pi(G_\la)$ over
$\la\vdash n$ grows (it is $5,9,11,15,15,21$ for $n=6,8,\dots,16$), attained at
``staircase'' $4$--cores such as $(3,2,1),(5,2,1),(5,2,1,1,1),(10,3,1,1,1)$.
Hence there is \emph{no} reduction to a finite size--bounded check, and the
termination statement is irreducibly uniform--in--$n$.

\paragraph{Gap B (two rows, uniform).}  Prove $r_b\ne r_{b-1}$ for all
$0\le b\le m$ with $(m,b)\ne(2,2)$, where $r_l=[u^l](u^2+(1+i)u+1)^m$.  This is the
two--row specialisation of Gap~A and is fully explicit but requires controlling the
simultaneous vanishing of $\mathrm{Re}$ and $\mathrm{Im}$ of a $\Z[i]$--binomial sum.

\paragraph{Why the $d=3$ method does not transplant.}  At $\zeta_3$ the analogous
value is a manifestly Schur--positive (hence nonzero) count, killing all
cancellation.  At $\zeta_4$ the value \eqref{eq:gauss} is genuinely complex, so no
positivity is available; the cancellation is real and must be controlled by a
valuation/Newton--polygon argument, which is precisely what Gap~A formalises.

\section{Computational verification}
All claims were verified with exact arithmetic (Pieri chain enumeration for $c_v$,
$M_j$; exact $\Z[i]$ for $G$):
\begin{itemize}
\item $(2,2)$ is the unique vanisher for all $\la\vdash n$, $n\le16$ ($913$ shapes).
\item Two--row formula of Prop.~\ref{prop:tworow} matches direct computation for all
$m\le10$; the only two--row vanisher up to $m=29$ is $(2,2)$.
\item Theorem~\ref{thm:partial}: $212/294$ shapes ($n\le14$) have $|J^\star|=1$;
all are nonzero with $v_\pi=V$.  Among the $82$ ties the only vanisher is $(2,2)$.
\end{itemize}

\section*{Summary}
The contribution of this note is a set of clean, mutually reinforcing
reformulations (the polynomial $P_\la$, the coproduct
$\Delta\psi=\psi\otimes1+1\otimes\psi+(1+i)p_1\otimes p_1$, the multivariate
generating function \eqref{eq:genfun}), a complete reduction and near--complete
proof of the two--row case, and an unconditional proof of non--vanishing for the
$72\%$ of shapes with a unique Newton vertex.  The biconditional reduces to the
single statement Gap~A: the second--order $(1+i)$--adic cancellation terminates for
every $\la\ne(2,2)$.
\end{document}
